%% 第 25 章：极分解
%% 本文件由 tools/md2tex.py 自动生成，请勿直接编辑。

\chapter{极分解}


\section{为什么需要极分解}

谱定理只处理：

\[
T:V\to V.
\]

但大量线性映射的定义域与陪域不同：

\[
T:V\to W.
\]

极分解将一般映射拆成两件事：

\begin{jitem}
\item %
输入方向被拉伸多少；
\item %
拉伸后如何以保长度的方式搬运到输出空间。
\end{jitem}

\begin{jdefn}{绝对值算子}{r:26:1}
对：

\[
T:V\to W,
\]

算子：

\[
T^*T:V\to V
\]

是正算子，因为：

\[
\langle T^*Tv,v\rangle
=
\langle Tv,Tv\rangle
=
\|Tv\|^2
\ge0.
\]

因此可定义：

\[
|T|
=
(T^*T)^{1/2}.
\]

它称为 $T$ 的绝对值算子。
\end{jdefn}

\begin{jthm}{极分解定理}{r:26:2}
存在唯一的部分等距算子：

\[
U:V\to W,
\]

使：

\[
T=U|T|.
\]

其中：

\[
|T|=(T^*T)^{1/2}.
\]
\end{jthm}

\section{极分解的意义}

$|T|$ 是定义在输入空间上的正算子，它描述各输入方向被放大的程度。

$U$ 则负责方向的搬运。

更准确地说，$U$ 在：

\[
(\ker T)^\perp
\]

到：

\[
\operatorname{range}T
\]

之间是等距同构；而在：

\[
\ker T
\]

上为零。

若 $T$ 可逆，则：

\[
U
\]

是整个空间上的酉算子或正交算子。
